% ===================================================================
%  GAMBAR No. 15 — Pasar. --- Panganan.
%
%  Traduction, faite en 2026, en javanais d'aujourd'hui, sur l'ido de
%  texto/io. Regles de la colonne : voir 10-gambar-01.tex. Une seule
%  numerotation.
%
%  « pasar » PARAIT ICI POUR LA QUATRIEME FOIS DANS LE RELEVE, ET LA
%  TROISIEME LANGUE A LE PORTER. Le persan bāzār est arrive en ourdou
%  par la terre, en indonesien et en javanais par la mer avec les
%  marchands venus de l'ouest, et en ido par l'italien. Le javanais le
%  tient de la meme route que l'indonesien, mais il l'a fait sien plus
%  loin : « pasar malem », la fete foraine du tableau 3, et
%  « pasaran », le cycle de cinq jours du calendrier javanais, dont
%  les noms — Legi, Pahing, Pon, Wage, Kliwon — sont ceux des marches
%  qui se tenaient chacun son jour. Un mot persan est devenu ici une
%  unite de temps.
%
%  ET LE POTAGER EST L'ENVERS EXACT DE CELUI DU TABLEAU 15
%  INDONESIEN. La-bas, wortel, buncis, kol, kembang kol, prei, seledri
%  etaient neerlandais d'un bout a l'autre, parce que ce sont les
%  Hollandais qui ont plante ces legumes dans les jardins de Java. Le
%  javanais dit exactement les memes mots — la couche est la meme, les
%  jardins etaient les memes —, et c'est la seule liste du livret ou
%  les deux colonnes coincident presque mot pour mot. Deux langues qui
%  ne partagent pas un mot sur trente parties du corps (tableau 2)
%  partagent presque tout le potager.
%
%  MAIS LE FOND JAVANAIS TIENT LA NOURRITURE : beras et sega le riz
%  (58), gabah, iwak le poisson, daging la viande, endhog les oeufs
%  (67), gula le sucre, uyah le sel, brambang et bawang l'oignon et
%  l'ail (101), lombok, tempe, tahu, jangan la soupe, panganan. Ce
%  qu'on mange tous les jours n'a jamais eu besoin d'un autre nom.
%
%  ET LE RIZ (58) EST LE MOT QUI SEPARE LE MIEUX CETTE COLONNE DE
%  TOUTES LES AUTRES DU RELEVE. Le javanais en a QUATRE — pari sur
%  pied, gabah en grain non decortique, beras cru et decortique, sega
%  cuit —, la ou l'ido, le francais et l'anglais n'en ont qu'un. La
%  planche grave un sac de riz d'epicerie : c'est « beras », et pas un
%  des trois autres. Le mot juste ne se choisit pas dans un
%  dictionnaire, il se choisit sur l'objet.
%
%  LES FRUITS EUROPEENS S'EMPRUNTENT TOUS — pir (52), plum (53),
%  anggur (54), persik (55), amandel (56), figa (64), zaitun (68) —
%  ET LA BANANE (46) EST LE SEUL FRUIT DU TABLEAU QUI SOIT D'ICI :
%  gedhang. L'ido la range parmi les « frukti exotika » ; pour cette
%  colonne, l'exotique est tout le reste de l'etal.
% ===================================================================

%%K t15-apar-1 apar
\VUsaut{4.0mm}
\VUtitre{15.0pt}{60}{GAMBAR No. 15}
\VUsaut{3.0mm}

%%K t15-tit-1 sub
\VUpk{11.4pt}{40}{Pasar. --- Panganan.}
\VUsaut{3.0mm}

%%K t15-01-1 p
1. --- Pedagang kang paling akeh, kang nyedhiyani kita barang kang
perlu kanggo pangan kita, padha nglumpuk ing sangisore \VUgras{aula}
\textsuperscript{(1)} \VUgras{pasar kang katutupan} utawa ing dalan-dalan
sacedhake.

%%K t15-02-1 p
2. --- \VUgras{Bakul daging babi} \textsuperscript{(2)}, kang dodol
\VUgras{ham} \textsuperscript{(3)}, \VUgras{sosis getih} \textsuperscript{(4)},
\VUgras{sosis cilik} \textsuperscript{(5)}, \VUgras{sosis usus}
\textsuperscript{(6)} lan \VUgras{lemak} \textsuperscript{(7)}, ngadeg
cedhak \VUgras{bakul mentega lan keju} \textsuperscript{(8)}, kang
pajangane kebak \VUgras{keju Walanda} \textsuperscript{(9)} kang kulite
abang, \VUgras{keju Camembert} \textsuperscript{(10)}, \VUgras{bunder
Gruyère} \textsuperscript{(11)} lan \VUgras{blok mentega}
\textsuperscript{(12)} kang seger.

%%K t15-03-1 p
3. --- Ing ngarepe wong iku, \VUgras{bakul sayur} \textsuperscript{(13)}
nawakake marang para tamune \VUgras{tomat} \textsuperscript{(14)} kang
endah, \VUgras{kembang kol} \textsuperscript{(15)}, \VUgras{salad}
\textsuperscript{(16)}, \VUgras{kol} \textsuperscript{(17)},
\VUgras{waluh} \textsuperscript{(19)}, \VUgras{semangka}
\textsuperscript{(18)} lan malah \VUgras{jamur} \textsuperscript{(20)}.
Nanging \VUgras{kancane tukang bir,} kang mandhegake \VUgras{kreta
kiriman}e \textsuperscript{(21)} kang kebak \VUgras{tong bir}
\textsuperscript{(22)} pas ing ngarepe pajangane, ngalang-alangi para
tamune nyedhak. Ana tukang kebon wadon nggawakake sakranjang
\VUgras{artisok} \textsuperscript{(23)} marang dheweke.

%%K t15-tit-2 sub
\VUsaut{4.0mm}
\VUpk{10.2pt}{40}{Dodol lan Tuku.}
\VUsaut{3.0mm}

%%K t15-04-1 p
4. --- Ing pasar, rame banget, amarga wis tekan jame \VUgras{dodolan
lelang} \textsuperscript{(24)}. \VUgras{Para nyonyah omah}
\textsuperscript{(26)}, kang teka mrana kanggo blanja, mandheg nalika
liwat ing ngarepe warunge \VUgras{bakul jeroan} \textsuperscript{(25)},
kanggo tuku \VUgras{babat} utawa \VUgras{sirah pedhet.}

%%K t15-05-1 p
5. --- \VUgras{Juru masak} \textsuperscript{(27)} kang lemu nawar
\VUgras{iwak tuna} kang endah banget marang \VUgras{bakul iwak}
\textsuperscript{(28)} kang sasenenge dhewe munggahake regane. Dheweke
nampa akeh banget \VUgras{welut, iwak sole, iwak trout} lan
\VUgras{iwak mas.} \VUgras{Iwak kod}e lan \VUgras{kerang}e seger banget.

%%K t15-tit-3 sub
\VUsaut{4.0mm}
\VUpk{10.2pt}{40}{Barang murah lan barang larang.}
\VUsaut{3.0mm}

%%K t15-06-1 p
6. --- \VUgras{Bakul unggas} \textsuperscript{(29)}, kang manggon ing
cedhake, jujur banget. Yen dheweke dodol \VUgras{kalkun, banyak,
bebek} utawa mung \VUgras{pitik enom,} dheweke njaluk rega kang pantes
wae.

%%K t15-07-1 p
7. --- Ing pojoke alun-alun, ing lantai kapisan, wis sawetara wektu
manggon \VUgras{bakul lawon} \textsuperscript{(30)}, panggonan para
wong tuku wadon tansah mesthi nemu pilihan kang becik banget bab
\VUgras{taplak meja, serbet, clemek} lan \VUgras{lap.} Ing lantai kang
padha, \VUgras{tukang ladhing} \textsuperscript{(31)} mbukak papan
gawene. Dheweke nglandhepake \VUgras{ladhing} para bakul wadon ing aula
lan gawe \VUgras{arit} saka \VUgras{waja} kang paling atos kanggo para
tukang kebon.

%%K t15-tit-4 sub
\VUsaut{4.0mm}
\VUpk{10.2pt}{40}{Toko bumbu.}
\VUsaut{3.0mm}

%%K t15-08-1 p
8. --- Ing sangisore papan gawene, wis sawetara wektu dibukak toko
panganan kang gedhe. Toko iku wis dudu \VUgras{toko bumbu}
\textsuperscript{(32)} kang prasaja lan ora wigati kaya biyen, kang mung
dodol barang lumrah — \VUgras{teh} \textsuperscript{(33)},
\VUgras{coklat} \textsuperscript{(34)}, \VUgras{gula batu}
\textsuperscript{(37)} lan \VUgras{sabun} \textsuperscript{(38)}. Ing kono
wong dodol barang apa wae : \VUgras{roti garing,}
\VUgras{panganan kaleng} \textsuperscript{(35)}, \VUgras{ciu}
\textsuperscript{(36)}, \VUgras{rum, konyak, kirsy, ciu adas} lan
\VUgras{kurasao.} Ing kono wong nemu buron : \VUgras{terwelu alas}
\textsuperscript{(39)}, \VUgras{manuk turdus} \textsuperscript{(40)},
\VUgras{ayam alas} \textsuperscript{(41)}, \VUgras{manuk gemak}
\textsuperscript{(42)}, \VUgras{bebek alas} \textsuperscript{(43)} utawa
\VUgras{manuk merak alas} \textsuperscript{(44)}, padha gampange karo
woh-wohan manca kayata \VUgras{gedhang} \textsuperscript{(46)} lan
\VUgras{nanas.}

%%K t15-09-1 p
9. --- \VUgras{Kancane} \textsuperscript{(45)} tukang bumbu kang grapyak
banget, siyap menehake apa kang dikarepake wong : \VUgras{selai}
\textsuperscript{(47)} saka woh anggur utawa woh quince, \VUgras{gajih}
\textsuperscript{(48)}, \VUgras{timun acar} \textsuperscript{(49)} utawa
\VUgras{sardin} \textsuperscript{(50)} asin.

%%K t15-09-2 p
Iku kepenak banget kanggo \VUgras{para tamu wadon} \textsuperscript{(51)},
kang nemu, ing sandhinge barang kang paling lumrah, panganan kang
paling arang lan woh-wohan kang paling enak : \VUgras{pir}
\textsuperscript{(52)} kang akeh banyune, \VUgras{plum}
\textsuperscript{(53)}, \VUgras{anggur} \textsuperscript{(54)},
\VUgras{persik} \textsuperscript{(55)}, \VUgras{amandel}
\textsuperscript{(56)} lan \VUgras{kenari} \textsuperscript{(57)}.

%%K t15-tit-5 sub
\VUsaut{4.0mm}
\VUpk{10.2pt}{40}{Para Tamu.}
\VUsaut{3.0mm}

%%K t15-10-1 p
10. --- \VUgras{Beras} \textsuperscript{(58)} lan \VUgras{kacang lentil}
\textsuperscript{(59)} ana ing sandhinge peti \VUgras{iwak haring}
\textsuperscript{(60)} kang diasapi ; \VUgras{kentang}
\textsuperscript{(61)} lan \VUgras{buncis} \textsuperscript{(63)} ana ing
cedhake \VUgras{apel} \textsuperscript{(62)}, \VUgras{woh figa}
\textsuperscript{(64)}, \VUgras{plum garing} \textsuperscript{(65)} lan
\VUgras{kacang hazel} \textsuperscript{(66)}. Ing toko iku wong nemu
\VUgras{endhog} \textsuperscript{(67)} anyar lan \VUgras{zaitun}
\textsuperscript{(68)} ; mula \VUgras{kranjang} \textsuperscript{(70)} lan
\VUgras{jaring blanja} \textsuperscript{(73)} enggal banget kebak.

%%K t15-11-1 p
11. --- \VUgras{Kopi} \textsuperscript{(72)} saka toko kang becik iku wis
oleh jeneng kang pantes ing antarane \VUgras{para batur wadon}
\textsuperscript{(69)} lan \VUgras{para juru masak wadon,} amarga
\VUgras{tukang bumbu} \textsuperscript{(71)} tuwa iku dhewe kang nggoreng
kopine kanthi ngati-ati banget.

%%K t15-tit-6 sub
\VUsaut{4.0mm}
\VUpk{10.2pt}{40}{Tontonan.}
\VUsaut{3.0mm}

%%K t15-12-1 p
12. --- Ora kabeh wong menyang pasar kanggo blanja. Ing lakune wong
kang endah ing dalan cilik kang tumuju aula, wong weruh wong kang
maneka warna banget. Ing sandhinge \VUgras{batur tukang daging}
\textsuperscript{(75)} kang grapyak, kang nyurung \VUgras{gerobag
dorong} \textsuperscript{(74)}, iki prajurit loro kang lagi prei lan
bangga banget nuduhake seragame kang mencorong : \VUgras{prajurit
husar} \textsuperscript{(76)} karo \VUgras{jas dolman} biru lan
\VUgras{topi cako} kang ana wulune, sarta \VUgras{kopral zuaf,} kang
kathoke mekrok lan sirahe katutupan \VUgras{kopyah fez.}

%%K t15-13-1 p
13. --- \VUgras{Tukang roti} \textsuperscript{(80)}, kang ngadeg ing
tlundhake \VUgras{toko roti} \textsuperscript{(78)}, lagi wae nguleni
jladren \VUgras{roti}ne \textsuperscript{(79)} lan njupuk saka oven
\VUgras{roti bulan sabit} \textsuperscript{(81)}, \VUgras{roti manis}
\textsuperscript{(82)} lan \VUgras{roti bunder} \textsuperscript{(83)} kang
saiki ana ing etalase.

%%K t15-14-1 p
14. --- Ing dhuwure toko roti manggon \VUgras{tukang jait lawon}
\textsuperscript{(84)} kang trampil, kang nyambut gawe karo sawetara
\VUgras{tukang bordhir wadon} \textsuperscript{(85)}. Ing sandhinge
manggon \VUgras{tukang umbah lan nyetrika} \textsuperscript{(86)}, kang
ngangkani \VUgras{sandhangan lawon} \textsuperscript{(88)} sawise diumbah
lan dipepe, lan kang nyetrika nganggo \VUgras{setrika}
\textsuperscript{(87)}.

%%K t15-tit-7 sub
\VUsaut{4.0mm}
\VUpk{10.2pt}{40}{Daging, iwak lan sayuran.}
\VUsaut{3.0mm}

%%K t15-15-1 p
15. --- Ing \VUgras{papan daging} \textsuperscript{(89)}, kang ana ing
lantai ngisor, wong dodol sakabehing \VUgras{daging} — \VUgras{daging
wedhus} \textsuperscript{(90)}, \VUgras{daging sapi} \textsuperscript{(94)}
utawa \VUgras{daging pedhet} \textsuperscript{(92)} — arupa \VUgras{pupu
wedhus, bistik, panggangan} lan \VUgras{kotelet.} \VUgras{Kancane
tukang daging} \textsuperscript{(93)} mbedhah kabeh daging iku lan
misahake \VUgras{balung sungsum} \textsuperscript{(96)}. Ing lawange papan
daging ana \VUgras{bakul tiram} \textsuperscript{(97)} kang dodol
\VUgras{tiram} \textsuperscript{(98)}. Dheweke mbukakake yen dijaluk.

%%K t15-16-1 p
16. --- Kabeh pedagang iku mbayar ijin utawa pajeg panggonan ; mula
\VUgras{polisi} \textsuperscript{(100)} tanpa welas mburu \VUgras{bakul
sayur mider wadon} \textsuperscript{(99)} kang mesakake, kang dadi
saingane merga ngider nggawa ing kreta ciliké \VUgras{wortel, lobak,
lobak putih, prei} utawa \VUgras{bongkokan asparagus, kacang polong
seger, bayem, jukut asam, selada banyu} lan \VUgras{seledri.}

%%K t15-tit-8 sub
\VUsaut{4.0mm}
\VUpk{10.2pt}{40}{Awas polisi!}
\VUsaut{3.0mm}

%%K t15-17-1 p
17. --- Bocah kang dodol \VUgras{bawang putih} \textsuperscript{(101)} lan
\VUgras{brambang} ngerti banget yen dheweke uga ora kena dodol barange
cedhak pasar. Dheweke mlayu, karo bebaya nyandhung kranjang
\VUgras{yuyu, keteng} lan \VUgras{urang} duweke \VUgras{bakul}
\textsuperscript{(102)} \VUgras{urang barong} lan \VUgras{lobster.}

%%K t15-18-1 p
18. --- Kabeh wong obah lan rame banget ; nanging iku ora ngalangi wong
krungu cetha swarane \VUgras{tukang kaca} \textsuperscript{(103)} kang
ngider, lan bengokane \VUgras{tukang areng} \textsuperscript{(104)} kang
lagi wae ninggal gerobage sedhela kanggo ngeterake \VUgras{karung
areng} \textsuperscript{(105)}.

\VUsaut{6.0mm}
\VUfilet{22mm}
